\documentclass[paper=a4,fontsize=10pt,DIV11,BCOR10mm]{scrartcl}


\usepackage[english]{babel}
\usepackage[utf8]{inputenc}
\usepackage[T1]{fontenc}
\usepackage{lmodern}

\usepackage[onehalfspacing]{setspace}
\setlength{\parindent}{0pt}

\usepackage{amssymb}
\usepackage{amsmath}
\usepackage{amsthm}
\usepackage{cancel}

\usepackage{graphicx}

\usepackage{url}
\usepackage{eurosym}


\usepackage{url}

\usepackage{placeins}


\newcommand{\abs}[1]{\left\lvert#1\right\rvert}
\newcommand{\norm}[1]{\left\lVert#1\right\rVert}





\titlehead{Technische Universität Berlin -- Fachgebiet Maschinelles Lernen\hfill \parbox[t]{2cm}{\includegraphics[width=2cm]{../TU_Logo_kurz_RGB_rot}}}

\begin{document}

\title{Machine Learning I - Assignment 8\\
\small{Technische Universität Berlin}}


\author{\hspace{5cm}\textbf{Gruppe SciComp}\hspace{5cm} \and
	\small{Christoph Conrads (315565)} \and
	\small{Antje Relitz (327289)} \and
	\small{Benjamin Pietrowicz (332542)} \and
	\small{Mitja Richter (324680)}
}

\date{WS 2013/2014}

\maketitle


\section{Discrete EM, coin tosses from two distributions}

Since we want to find the parameter $\theta=(\lambda,p_1,p_2)$ that maximizes \begin{equation*}
Q(\theta,\theta_{old})=\sum_{z\in \{T,H\}^N}\underbrace{P(Z=z|X=x,\theta_{old})}_{=:q(z)}  \underbrace{\log(P(X=x,Z=z|\theta))}_{=:F(z,\theta)},
\end{equation*}
all we have to do is to rewrite $Q$ in terms of $\theta$ and then set the derivatives of $Q$ with respect to each component of $\theta$ to zero. We first note, that the factor $P(Z=z|X=x,\theta_{old})$ does not depend on $\theta$ and will therefore remain unchanged under differentiation. We thus only have to care about the function $F(z,\theta)$ for now.

We introduce the notation $H(v)\equiv$ number of heads in $v$ and $T(v)\equiv$ number of tails in $v$ for any vector $v$ (including scalars).

We do now obtain:
\begin{align*}
F(z,\theta)&=\log(P(X=x,Z=z|\theta))\\
&=\log( \prod_{i=1}^N P(Z=z^{(i)}|\theta) \prod_{j=1}^m P(X_j=x_j^{(i)}|Z=z^{(i)},\theta)) \\
&=\sum_{i=1}^N \log (P(Z=z^{(i)}|\theta)) + \sum_{j=1}^m \log (P(X_j=x_j^{(i)}|Z=z^{(i)},\theta)).
\end{align*}

With our prior knowledge, we can now rewrite
\begin{align*}
&P(Z=z^{(i)}|\theta)=\lambda^{H(z^{(i)})}(1-\lambda)^{T(z^{(i)})}\\
\text{and } & P(X_j=x_j^{(i)}|Z=z^{(i)},\theta)= \bigl( p_1^{H(x_j^{(i)})} (1-p_1)^{T(x_j^{(i)})}\bigr)^{H(z^{(i)})} \bigl(  p_2^{H(x_j^{(i)})} (1-p_2)^{T(x_j^{(i)})} \bigr)^{T(z^{(i)})}.
\end{align*}

If we insert this into the equation above, we get:

\begin{align*}
F(z,\theta)&=\sum_{i=1}^N \log (\lambda^{H(z^{(i)})}(1-\lambda)^{T(z^{(i)})})\\
& + \sum_{j=1}^m \log \Bigl(\bigl( p_1^{H(x_j^{(i)})} (1-p_1)^{T(x_j^{(i)})}\bigr)^{H(z^{(i)})} \bigl(  p_2^{H(x_j^{(i)})} (1-p_2)^{T(x_j^{(i)})} \bigr)^{T(z^{(i)})} \Bigr)\\
&=\sum_{i=1}^N H(z^{(i)})\log (\lambda)+ T(z^{(i)}) \log (1-\lambda) +\sum_{j=1}^m  H(x_j^{(i)}) H(z^{(i)}) \log(p_1)\\
& + T(x_j^{(i)}) H(z^{(i)})\log(1-p_1) + H(x_j^{(i)})T(z^{(i)}) \log(p_2)+ T(x_j^{(i)})T(z^{(i)}) \log (1-p_2).
\end{align*}

\subsection*{Determination of $\hat \lambda$}
We now want to derive an expression for $\hat \lambda$ and therefore we set the derivative of $Q$ w.r.t. $\lambda$ to zero:
\begin{equation}\label{derQlambda}
\frac{dQ(\theta,\theta_{old})}{d\lambda}=\sum_{z\in \{T,H\}^N} q(z) \frac{dF(z,\theta)}{d\lambda} \stackrel{!}{=}0.
\end{equation}
So we see that we have to compute the derivative of $F(z,\theta)$ w.r.t. $\lambda$.

\begin{align*}
\frac{dF(z,\theta)}{d\lambda} &=\frac{d}{d\lambda}\sum_{i=1}^N H(z^{(i)})\log (\lambda)+ T(z^{(i)}) \log (1-\lambda)\\
&=\frac{d}{d\lambda}\log (\lambda)\sum_{i=1}^N H(z^{(i)})+  \frac{d}{d\lambda}\log (1-\lambda)\sum_{i=1}^N T(z^{(i)})\\
&= \frac{H(z)}{\lambda} - \frac{T(z)}{1-\lambda}\\
&=\frac{(1-\lambda)H(z)-\lambda T(z)}{\lambda(1-\lambda)}\\
& \stackrel{T(z)=N-H(z)}{=}\frac{H(z)-\lambda N}{\lambda(1-\lambda)}
\end{align*}

We can now insert this result into equation \ref{derQlambda}:
\begin{align*}
\sum_{z\in \{T,H\}^N} q(z) \frac{dF(z,\theta)}{d\lambda} &=0\\
\Leftrightarrow \sum_{z\in \{T,H\}^N} q(z) \frac{H(z)-\lambda N}{\lambda(1-\lambda)} &=0\\
\Leftrightarrow \sum_{z\in \{T,H\}^N} q(z) \frac{H(z)}{\lambda(1-\lambda)} &= \sum_{z\in \{T,H\}^N} q(z) \frac{\lambda N}{\lambda(1-\lambda)}\\
\Leftrightarrow \sum_{z\in \{T,H\}^N} q(z) H(z) &= \lambda N \sum_{z\in \{T,H\}^N} q(z) \\
\Leftrightarrow \frac{\sum_{z\in \{T,H\}^N} q(z) H(z)}{N \sum_{z\in \{T,H\}^N} q(z)} &= \lambda\\
\end{align*}
Since $q(z)$ represents a probability distribution, we know that $\sum_{z\in \{T,H\}^N} q(z)=1$ and therefore we obtain
\begin{equation*}
\hat \lambda = \frac{\sum_{z\in \{T,H\}^N} q(z) H(z)}{N}.
\end{equation*}

\subsection*{Determination of $\hat p_1$}

We now want to derive an expression for $\hat p_1$ and therefore we set the derivative of $Q$ w.r.t. $p_1$ to zero:
\begin{equation}\label{derQp1}
\frac{dQ(\theta,\theta_{old})}{d p_1}=\sum_{z\in \{T,H\}^N} q(z) \frac{dF(z,\theta)}{d p_1} \stackrel{!}{=}0.
\end{equation}
So we see that we have to compute the derivative of $F(z,\theta)$ w.r.t. $p_1$.

\begin{align*}
\frac{dF(z,\theta)}{d p_1} &= \frac{d}{d p_1} \Bigl ( \sum_{i=1}^N \sum_{j=1}^m  H(x_j^{(i)}) H(z^{(i)}) \log(p_1) + T(x_j^{(i)}) H(z^{(i)})\log(1-p_1) \Bigr )\\
&= \frac{d}{d p_1} \log(p_1) \sum_{i=1}^N  H(z^{(i)})\sum_{j=1}^m  H(x_j^{(i)})  + \frac{d}{d p_1} \log(1-p_1) \sum_{i=1}^N H(z^{(i)}) \sum_{j=1}^m T(x_j^{(i)}) \\
&= \frac{1}{p_1} \sum_{i=1}^N  H(z^{(i)}) H(x^{(i)}) - \frac{1}{1-p_1} \sum_{i=1}^N H(z^{(i)}) T(x^{(i)}) \\
&\stackrel{T(x^{(i)})=m-H(x^{(i)})}{=} \frac{1}{p_1} \sum_{i=1}^N  H(z^{(i)}) H(x^{(i)}) - \frac{1}{1-p_1} \sum_{i=1}^N m H(z^{(i)}) + \frac{1}{1-p_1} \sum_{i=1}^N H(x^{(i)})H(z^{(i)}) \\
&=\frac{1}{p_1(1-p_1)} \sum_{i=1}^N  H(z^{(i)}) H(x^{(i)}) - \frac{m}{1-p_1} H(z).
\end{align*}

We can now insert this expression into equation \ref{derQp1}:
\begin{align*}
\sum_{z\in \{T,H\}^N} q(z) \frac{dF(z,\theta)}{d p_1} &=0\\
\Leftrightarrow \sum_{z\in \{T,H\}^N} q(z)\Bigl ( \frac{1}{p_1(1-p_1)} \sum_{i=1}^N  H(z^{(i)}) H(x^{(i)}) - \frac{m}{1-p_1} H(z) \Bigr) &=0\\
\Leftrightarrow \sum_{z\in \{T,H\}^N} q(z) \frac{1}{p_1(1-p_1)} \sum_{i=1}^N  H(z^{(i)}) H(x^{(i)}) &=\sum_{z\in \{T,H\}^N} q(z) \frac{m}{1-p_1} H(z)\\
\Leftrightarrow \sum_{z\in \{T,H\}^N} q(z) \sum_{i=1}^N  H(z^{(i)}) H(x^{(i)}) &=m \cdot p_1 \sum_{z\in \{T,H\}^N} q(z) H(z)
\end{align*}
 
 and therefore we get:
 \begin{align*}
 \hat p_1 &=\frac{ \sum_{z\in \{T,H\}^N} q(z) \sum_{i=1}^N  H(z^{(i)}) H(x^{(i)})}{m\sum_{z\in \{T,H\}^N} q(z) H(z)}\\
 &=\frac{ \sum_{z\in \{T,H\}^N} q(z) \sum_{i=1}^N  H(z^{(i)}) H(x^{(i)})}{m \cdot N  \cdot \hat \lambda}.
 \end{align*}

\subsection*{Determination of $\hat p_2$}

In order to determine an explicit expression for $\hat p_2$, we proceed analogously to the case of $\hat p_1$.

We set the derivative of $Q$ w.r.t. $p_2$ to zero:
\begin{equation}\label{derQp2}
\frac{dQ(\theta,\theta_{old})}{d p_2}=\sum_{z\in \{T,H\}^N} q(z) \frac{dF(z,\theta)}{d p_2} \stackrel{!}{=}0
\end{equation}
and compute the derivative of $F(z,\theta)$ w.r.t. $p_2$:

\begin{align*}
\frac{dF(z,\theta)}{d p_2} &= \frac{d}{d p_2} \Bigl ( \sum_{i=1}^N \sum_{j=1}^m  H(x_j^{(i)})T(z^{(i)}) \log(p_2)+ T(x_j^{(i)})T(z^{(i)}) \log (1-p_2) \Bigr )\\
&= \frac{d}{d p_2} \log(p_2)\sum_{i=1}^N \sum_{j=1}^m  H(x_j^{(i)})T(z^{(i)}) +  \frac{d}{d p_2} \log (1-p_2) \sum_{i=1}^N \sum_{j=1}^m T(x_j^{(i)})T(z^{(i)})  \\
&= \frac{1}{p_2}\sum_{i=1}^N  H(x^{(i)})T(z^{(i)}) -  \frac{1}{1- p_2}\sum_{i=1}^N  T(x^{(i)})T(z^{(i)})  \\
&\stackrel{T(x^{(i)})=m-H(x^{(i)})}{=} \frac{1}{p_2}\sum_{i=1}^N  H(x^{(i)})T(z^{(i)}) -  \frac{1}{1- p_2}\sum_{i=1}^N  m T(z^{(i)}) + \frac{1}{1- p_2}\sum_{i=1}^N H(x^{(i)})T(z^{(i)}) \\
&= \frac{1}{p_2(1-p_2)}\sum_{i=1}^N  H(x^{(i)})T(z^{(i)}) -  \frac{m}{1- p_2} T(z).
\end{align*}

We can now insert this into equation \ref{derQp2} and obtain:

\begin{align*}
\sum_{z\in \{T,H\}^N} q(z) \frac{dF(z,\theta)}{d p_2} &=0\\
\Leftrightarrow \sum_{z\in \{T,H\}^N} q(z)\Bigl ( \frac{1}{p_2(1-p_2)}\sum_{i=1}^N  H(x^{(i)})T(z^{(i)}) -  \frac{m}{1- p_2} T(z) \Bigr) &=0\\
\Leftrightarrow  \sum_{z\in \{T,H\}^N} q(z)\frac{1}{p_2(1-p_2)}\sum_{i=1}^N  H(x^{(i)})T(z^{(i)}) &=  \sum_{z\in \{T,H\}^N} q(z) \frac{m}{1- p_2} T(z)\\
\Leftrightarrow  \sum_{z\in \{T,H\}^N} q(z)\sum_{i=1}^N  H(x^{(i)})T(z^{(i)}) &=  m \cdot p_2\sum_{z\in \{T,H\}^N} q(z) T(z).
\end{align*}

This yields:

\begin{equation*}
\hat p_2 = \frac{\sum_{z\in \{T,H\}^N} q(z)\sum_{i=1}^N  H(x^{(i)})T(z^{(i)})}{m \sum_{z\in \{T,H\}^N} q(z) T(z)}
\end{equation*}

%\begin{thebibliograpy}
%	\bibitem{DHS2000} R. O. Duda, P. E. Hart and D. G. Stork. \emph{Pattern Classification}. 2nd ed. New York, NY, USA: Wiley-Interscience, 2000. ISBN: 0-4710-5669-3
%\end{thebibliograpy}

\end{document}
